\documentclass[12pt]{article}
\usepackage[utf8]{inputenc}
\usepackage[T1]{fontenc}
\usepackage{hyperref}
\title{Tiyaa Language Resources for K-12 Curriculum Development}
\author{Compiled by Claude Sonnet 5 (claude-sonnet-5) for Andr\'as Kornai}
\date{}
\begin{document}
\maketitle

\section{Introduction}
Tiyaa (ISO 639-3 tyy, Glottocode yaaa1242) is a Benue-Congo language of Nigeria. Glottolog places it in the Tarokoid branch of Plateau. The search results do not link it to the Kambari languages. Speaker estimates in the results are about 6,000 (ASJP) and about 7,100 (Joshua Project). Alternate names include Tiya, Ayaa and Ya. Very little material turned up.

\section{Primary Educational Resources}
I found no primers, textbooks or school programmes for Tiyaa.

\section{Linguistic and Reference Materials}
\begin{itemize}
\item ASJP word-list entry: \url{https://asjp.clld.org/languages/TIYAA}
\item Endangered Languages Project entry: \url{https://www.endangeredlanguages.com/bn/node/96036}
\end{itemize}

\section{Religious and Community Resources}
Joshua Project language profile: \url{https://joshuaproject.net/index.php/languages/tyy}

\section{Digital and Audio Resources}
I found none.

\section{Curriculum Development Context}
The documentation is minimal. Any mother-tongue curriculum would need basic description and orthography work first.
\end{document}